\section{Bilan de liaison et calcul de portée théorique}

L'argument principal en faveur du Wi-Fi HaLow (IEEE 802.11ah) pour le pilotage de drones FPV réside 
dans sa capacité à maintenir une liaison stable sur de grandes distances, y compris en présence 
d'obstacles. Pour valider mathématiquement cette capacité, un bilan de liaison (\textit{Link Budget}) a 
été calculé.

Ce calcul permet d'estimer la portée théorique maximale du système en condition de ligne de vue directe 
(\textit{Line-Of-Sight} - LOS), avant d'évaluer la marge disponible pour la pénétration des obstacles.

\subsection{Paramètres du système (Wi-Fi HaLow 915 MHz)}
Les performances RF de l'émetteur (MM8108) et du récepteur déterminent la puissance du signal. Les valeurs 
suivantes sont basées sur les spécifications techniques du matériel utilisé et la configuration logicielle 
retenue (MCS 2 / Bande passante de 2 MHz) :

\begin{itemize}
    \item \textbf{Fréquence d'opération ($f$) :} 915.5 MHz (Domaine US, Canal 27).
    \item \textbf{Puissance d'émission ($P_{tx}$) :} +20 dBm (100 mW), puissance typique maximale configurée 
    par le SDK.
    \item \textbf{Gain de l'antenne d'émission ($G_{tx}$) :} +2 dBi (Antenne dipôle standard).
    \item \textbf{Gain de l'antenne de réception ($G_{rx}$) :} +2 dBi.
    \item \textbf{Sensibilité du récepteur ($S_{rx}$) :} -98 dBm (Seuil minimum théorique pour décoder du 
    QPSK / MCS 2 sur 2 MHz).
    \item \textbf{Marge d'évanouissement (\textit{Fade Margin} - $F_m$) :} 10 dB. Cette marge de sécurité 
    est vitale dans les communications mobiles (drone) pour compenser les variations soudaines du signal dues 
    aux réflexions (trajets multiples) ou aux inclinaisons d'antennes.
\end{itemize}

\subsection{Calcul de l'atténuation maximale tolérable}
L'atténuation maximale tolérable par le système (le \textit{Maximum Path Loss}) correspond à la différence 
entre la puissance de sortie rayonnée et la limite de sensibilité du récepteur, diminuée de la marge de sécurité.

L'équation du bilan de liaison s'écrit :
$$L_{max} = P_{tx} + G_{tx} + G_{rx} - S_{rx} - F_m$$

En injectant nos paramètres matériels :
$$L_{max} = 20\text{ dBm} + 2\text{ dBi} + 2\text{ dBi} - (-98\text{ dBm}) - 10\text{ dB}$$
$$L_{max} = 112\text{ dB}$$

Le système peut donc subir une perte de signal de \textbf{112 dB} dans les airs avant que la liaison vidéo ne 
coupe de manière irrémédiable.

\subsection{Estimation de la portée en espace libre (FSPL)}
L'atténuation du signal dans l'air, sans aucun obstacle, est définie par le modèle de perte en espace libre 
(\textit{Free Space Path Loss} - FSPL). L'équation logarithmique du FSPL est la suivante :
$$FSPL(dB) = 20 \log_{10}(d) + 20 \log_{10}(f) + 32.44$$
Où :
\begin{itemize}
    \item $d$ est la distance en kilomètres (km).
    \item $f$ est la fréquence en mégahertz (MHz).
\end{itemize}

Pour trouver la distance maximale théorique ($d_{max}$), on pose l'équation $FSPL = L_{max}$ :
$$112 = 20 \log_{10}(d_{max}) + 20 \log_{10}(915.5) + 32.44$$
$$112 = 20 \log_{10}(d_{max}) + 59.23 + 32.44$$
$$112 = 20 \log_{10}(d_{max}) + 91.67$$
$$20 \log_{10}(d_{max}) = 20.33$$
$$\log_{10}(d_{max}) \approx 1.0165$$
$$d_{max} = 10^{1.0165} \approx \mathbf{10.38 \text{ kilomètres}}$$

\subsection{Validation empirique en environnement complexe (NLOS)}
Pour confronter la marge d'atténuation théorique à la réalité du terrain, des tests de portée ont été réalisés en 
conditions strictement non-optimales (\textit{Non-Line-Of-Sight} - NLOS) :
\begin{itemize}
    \item \textbf{Test 1 - Pénétration extrême (30 mètres) :} L'émetteur a été placé dans une chambre et le récepteur 
    dans une cave (mur en béton armé type "bunker"). La liaison vidéo s'est maintenue de manière fluide, bien qu'une 
    orientation de l'antenne vers les ouvertures ait été nécessaire pour stabiliser le flux. Ce test prouve la capacité 
    du 915 MHz à traverser des dalles de béton.
    \item \textbf{Test 2 - Obstacles multiples (160 mètres) :} L'émetteur est resté à l'intérieur (derrière un mur de ferme 
    de 40 cm d'épaisseur) tandis que le récepteur a été éloigné dans un champ à 160 mètres, à travers des arbres et soumis 
    aux réflexions (trajets multiples) des maisons environnantes. L'image est restée parfaitement fluide jusqu'à la limite 
    physique où la coupure est intervenue brutalement (comportement typique d'une démodulation numérique arrivant au seuil 
    de sa sensibilité).
\end{itemize}
Ces essais confirment que l'immense surplus de puissance (budget de liaison) calculé précédemment est effectivement "consommé" 
par la densité des obstacles, permettant une liaison vidéo FPV parfaitement utilisable dans des milieux où les ondes 5.8 GHz 
classiques ne passeraient même pas le premier mur.

\subsection{Conclusion du bilan de liaison}
Le calcul théorique démontre que, sous des conditions idéales de ligne de vue directe (LOS) avec une marge de 
sécurité de 10 dB, le système vidéo FPV basé sur le Wi-Fi HaLow peut porter jusqu'à \textbf{10 kilomètres}.

En pratique, pour le pilotage d'un drone en environnement complexe (arbres, bâtiments, murs d'une ferme), 
l'opérateur n'atteindra pas ces 10 km. Cependant, cette immense réserve de puissance (\textit{Link Budget} 
excédentaire) sera absorbée par les obstacles. C'est précisément ce surplus de puissance, couplé à la longueur 
d'onde pénétrante du Sub-1 GHz, qui garantit la stabilité du flux vidéo en situation de NLOS 
(\textit{Non-Line-Of-Sight}) à des distances de plusieurs centaines de mètres, là où les fréquences 
5.8 GHz standard échouent totalement.